\chapter{Relaciones de Recurrencia}

\section{Definiciones}

\begin{definicion}
Una \textbf{relación de recurrencia} para la sucesión $a_0, a_1, a_2, \ldots$ es una ecuación que relaciona $a_n$ con predecesores. Las \textbf{condiciones iniciales} son valores dados explícitamente para un número finito de términos.
\end{definicion}

\begin{ejemplo}[Fibonacci]
\[
f_n = f_{n-1} + f_{n-2}, \quad n \geq 3, \quad f_1 = f_2 = 1.
\]
\end{ejemplo}

\begin{ejemplo}[Torres de Hanoi]
\[
h_n = 2h_{n-1} + 1, \quad h_1 = 1.
\]
\end{ejemplo}

\section{Solución de Recurrencias}

\begin{definicion}
Resolver una relación de recurrencia significa hallar una fórmula explícita para $a_n$.
\end{definicion}

\begin{ejemplo}
$a_n = a_{n-1} + 3$, $a_1 = 2$ $\Rightarrow$ $a_n = 2 + 3(n-1)$.
\end{ejemplo}

\begin{ejemplo}
$h_n = 2h_{n-1} + 1$, $h_1 = 1$ $\Rightarrow$ $h_n = 2^n - 1$.
\end{ejemplo}

\begin{ejemplo}
$a_n = 3n\,a_{n-1}$, $a_0 = 2$ $\Rightarrow$ $a_n = 2 \cdot 3^n n!$.
\end{ejemplo}

\section{Recurrencias Homogéneas Lineales}

\begin{definicion}
Una recurrencia homogénea lineal de orden $k$ con coeficientes constantes es de la forma
\[
a_n = c_1 a_{n-1} + \cdots + c_k a_{n-k}, \quad c_k \neq 0.
\]
\end{definicion}

\begin{teorema}
Para la recurrencia de orden 2, $a_n = c_1 a_{n-1} + c_2 a_{n-2}$:
\begin{enumerate}
  \item Si $U_n, V_n$ son soluciones, $bU_n + dV_n$ también lo es.
  \item Si $r$ es raíz de $t^2 - c_1 t - c_2 = 0$, entonces $r^n$ es solución.
  \item Si $r_1 \neq r_2$ reales: $a_n = b r_1^n + d r_2^n$.
  \item Si $r$ es raíz doble: $a_n = b r^n + d n r^n$.
  \item Si las raíces son complejas: solución trigonométrica.
\end{enumerate}
\end{teorema}

\begin{ejemplo}[Fibonacci]
$f_n = f_{n-1} + f_{n-2}$, $f_0 = 0$, $f_1 = 1$. La ecuación característica es $r^2 - r - 1 = 0$, cuyas raíces son $\frac{1 \pm \sqrt{5}}{2}$. Por tanto,
\[
f_n = \frac{1}{\sqrt{5}}\left(\frac{1+\sqrt{5}}{2}\right)^n - \frac{1}{\sqrt{5}}\left(\frac{1-\sqrt{5}}{2}\right)^n.
\]
\end{ejemplo}

\begin{ejemplo}
$a_n = 6a_{n-1} - 9a_{n-2}$, $a_0 = 1$, $a_1 = -1$. La ecuación característica da $r = 3$ doble. Así $a_n = (b + dn)3^n$. Con las condiciones iniciales, $b = 1$, $d = -4/3$. Por tanto,
\[
a_n = \left(1 - \tfrac{4}{3}n\right)3^n.
\]
\end{ejemplo}

\section{Recurrencias No Homogéneas}

\begin{teorema}
Dada $a_n = c_1 a_{n-1} + c_2 a_{n-2} + f(n)$, si $g(n)$ es solución particular y $V_n$ es solución de la homogénea asociada, entonces $U_n = V_n + g(n)$ es solución general.
\end{teorema}

\begin{ejemplo}
$a_n = 6a_{n-1} - 8a_{n-2} + 3$, $a_0 = 0$, $a_1 = 1$.
Homogénea: $r^2 - 6r + 8 = 0 \Rightarrow r = 2, 4$.
Solución particular constante: $A = 6A - 8A + 3 \Rightarrow A = -1$.
Con las condiciones iniciales, $b = 0$, $d = 1/2$. Luego,
\[
a_n = \tfrac{1}{2}\,4^n - 1.
\]
\end{ejemplo}